\chapter{A Concrete Naming Certificate for $\FiniteIntervalMeshRefinementUp$}
\label{ch:concrete-instances-finite-interval-mesh-refinement-namecert}
\origin{human}

Following Bishop and Bridges, finite interval arguments are read through displayed meshes, endpoint comparisons, and rational tolerances rather than through a completed continuum of hidden subdivision choices. The $\FiniteIntervalMeshRefinementUp$ packet sits between the dyadic arithmetic of \autoref{ch:concrete-instances-dyadicratcore-namecert}, the compact interval surfaces of \autoref{ch:concrete-instances-compact-interval-modulus-namecert}, the compact metric finite-net route of \autoref{ch:concrete-instances-compactmetric-namecert}, the total-boundedness surface of \autoref{ch:concrete-instances-totallybounded-namecert}, and the terminal real-name seal of \autoref{ch:concrete-instances-real-namecert}. It names the finite refinement discipline used by compact interval and dyadic-net consumers without asserting open-cover compactness, selected partitions, or ambient $\RealUp$ completeness.

\begin{definition}[Finite interval mesh refinement packet]
\label{def:finite-interval-mesh-refinement-packet}
A $\FiniteIntervalMeshRefinementUp$ carrier is a finite $\BHist$ packet
$$
M=(I,D,Q,R,U,H,C,P,N)
$$
with the following displayed rows:
\begin{enumerate}
\item $I$ is the $\CompactIntervalUp$ endpoint row, displaying the interval source and the endpoint order data used by the mesh read.
\item $D$ is the $\DyadicRatCoreUp$ mesh ledger, listing the dyadic breakpoints, dyadic cell radii, endpoint incidence, and finite containment comparisons.
\item $Q$ is the refinement map from a coarse finite mesh to the displayed finer mesh, including the finite parent-cell table.
\item $R$ is the stability row saying that each refined dyadic cell inherits its endpoint, radius, and incidence rows from the parent cell through the same dyadic comparisons.
\item $U$ is the coverage row, recording that the refined mesh covers the same compact interval source and that a downstream $\TotallyBoundedUp$ or $\CompactMetricUp$ consumer reads only the displayed finite cells.
\item $H$ records componentwise $\hsame$ transport for interval endpoints, dyadic mesh rows, refinement maps, stability rows, coverage rows, replay, provenance, and local naming.
\item $C$ records displayed $\Cont$ replay from an interval-mesh request through $I,D,Q,R,U,H$.
\item $P$ is the $\Pkg$ provenance over $\FiniteIntervalMeshRefinementUp$, $\CompactIntervalUp$, $\DyadicRatCoreUp$, $\CompactMetricUp$, $\TotallyBoundedUp$, $\RealUp$, $\BHist$, $\hsame$, $\Cont$, and $\NameCert$ rows.
\item $N$ is the local $\NameCert_{\FiniteIntervalMeshRefinementUp}$ row.
\end{enumerate}
The carrier has no coordinate for an infinite subdivision, a selected limiting mesh, an open-cover choice principle, quotient equality of intervals, or ambient $\RealUp$ completeness.
\end{definition}

\begin{theorem}[Finite interval mesh refinement NameCert obligations]
\label{thm:finite-interval-mesh-refinement-namecert-obligations}
For every accepted $\FiniteIntervalMeshRefinementUp$ carrier
$$
M=(I,D,Q,R,U,H,C,P,N),
$$
the local $\NameCert_{\FiniteIntervalMeshRefinementUp}$ surface consists exactly of five finite obligation rows:
\begin{enumerate}
\item carrier admission for the compact interval source $I$, dyadic mesh ledger $D$, refinement table $Q$, stability row $R$, coverage row $U$, transport row $H$, replay row $C$, provenance row $P$, and local name $N$;
\item mesh refinement stability, requiring every refined cell in $Q$ to inherit its dyadic endpoint, radius, and incidence rows from a displayed parent cell through $R$;
\item coverage monotonicity, requiring the refined mesh to cover the same interval source as the coarse mesh and to expose no smaller domain when read by a $\CompactMetricUp$ or $\TotallyBoundedUp$ consumer;
\item $\hsame$ and $\Cont$ transport, requiring $H$ and $C$ to preserve the same interval source, dyadic rows, refinement table, stability row, coverage row, provenance row, and local name under accepted renamings and replay;
\item ledger refusal, requiring the packet to export no infinite subdivision, selected limiting mesh, open-cover choice principle, quotient interval equality, hidden Real selector, or ambient $\RealUp$ completeness.
\end{enumerate}
No finite interval mesh-refinement obligation is stored outside these rows.
\end{theorem}
\begin{proof}
Project the carrier of \autoref{def:finite-interval-mesh-refinement-packet}. The interval row $I$, dyadic ledger $D$, refinement table $Q$, stability row $R$, coverage row $U$, and structural rows $H,C,P,N$ are the entire displayed packet.

Carrier admission is therefore admission of those rows. Mesh refinement stability is read from $Q$ together with $R$: each refined cell points to a displayed parent cell, and the endpoint, radius, and incidence data are transported through dyadic comparisons already present in $D$.

Coverage monotonicity is read from $U$. The refined mesh may subdivide cells, but it keeps the same compact interval source and the same finite coverage target before any compact metric or totally bounded consumer reads it.

Finally $H,C,P,N$ transport, replay, source, and locally name the same finite route. Since the carrier contains no coordinate for infinite subdivision, selected limiting meshes, open-cover choice, quotient interval equality, hidden Real selectors, or ambient Real completeness, the five listed rows exhaust the local NameCert surface.
\end{proof}

\begin{theorem}[Finite interval mesh refinement compact handoff]
\label{thm:finite-interval-mesh-refinement-compact-handoff}
Every compact interval, compact metric, totally bounded, dyadic-net, or $\RealUp$-facing consumer of an accepted $\FiniteIntervalMeshRefinementUp$ packet must read the displayed route
$$
I,D,Q,R,U,H,C,P,N
$$
before treating a refined mesh as usable finite data. The handoff is a finite BEDC mesh route, not a selected subdivision limit, not an open-cover compactness theorem, and not an appeal to ambient $\RealUp$ completeness.
\end{theorem}
\begin{proof}
Start with the NameCert obligations of \autoref{thm:finite-interval-mesh-refinement-namecert-obligations}. They force the refined mesh to expose carrier admission, refinement stability, coverage monotonicity, transport, replay, provenance, local naming, and ledger refusal.

A compact interval or dyadic-net consumer reads $I,D,Q,R,U$ first. These rows display the source interval, dyadic cells, parent-cell table, inherited endpoint data, and coverage ledger.

A compact metric, totally bounded, or Real-facing consumer is downstream of the same displayed route. The rows $H,C,P,N$ preserve and name it without creating a selected limiting partition or a hidden completeness principle. Thus every accepted handoff remains finite and local to the packet.
\end{proof}

\falsifiablePrediction{Every accepted $\FiniteIntervalMeshRefinementUp$ read exposes a compact interval source row, dyadic mesh ledger, finite refinement table, stability row, coverage row, hsame transport, Cont replay, provenance, and local NameCert before a compact interval or Real-facing consumer reads the refined mesh.}

\independenceWitness{dyadicratcore, compact-interval, compactmetric, totallybounded, real: no listed sibling simultaneously carries a finite parent-cell refinement table, inherited dyadic endpoint rows, coverage monotonicity, transport, replay, provenance, and local NameCert rows for compact interval mesh refinement.}

\closureat{FiniteIntervalMeshRefinementUp}{seedStr}

\begin{closurestatus}{\FiniteIntervalMeshRefinementUp}
  \constructivestory{A $\FiniteIntervalMeshRefinementUp$ packet is generated from a compact interval source, a DyadicRatCore mesh ledger, a finite parent-cell refinement table, stability and coverage rows, componentwise hsame transport, Cont replay, Pkg provenance, and a local NameCert row.}
  \theoryclosure{\seedClosure}
  \scopeclosed{\autoref{def:finite-interval-mesh-refinement-packet}, \autoref{thm:finite-interval-mesh-refinement-namecert-obligations}, and \autoref{thm:finite-interval-mesh-refinement-compact-handoff} seed the finite compact-interval mesh refinement route from dyadic cell rows to compact and Real-facing consumers.}
  \formalstatus{\unformalizedV}
  \bridgestatus{none}
  \notclaimed{This chapter does not claim open-cover compactness, infinite subdivisions, selected limiting meshes, quotient interval equality, ambient RealUp completeness, Lean verification, or bridge checking.}
  \upgradepath{Obligation closure requires checked carrier admission, mesh refinement stability, coverage monotonicity, transport, replay, provenance, and ledger refusal for BEDC.Derived.FiniteIntervalMeshRefinementUp.}
\end{closurestatus}
